\documentclass[11pt]{article}

% ---------- Packages ----------
\usepackage[a4paper,margin=1in]{geometry}
\usepackage{amsmath,amssymb,amsthm,mathtools}
\usepackage{hyperref}
\usepackage{enumitem}
\usepackage[numbers]{natbib}

% ---------- Theorem environments ----------
\newtheorem{theorem}{Theorem}[section]
\newtheorem{proposition}[theorem]{Proposition}
\newtheorem{lemma}[theorem]{Lemma}
\newtheorem{corollary}[theorem]{Corollary}

\theoremstyle{definition}
\newtheorem{definition}[theorem]{Definition}
\newtheorem{remark}[theorem]{Remark}

% ---------- Commands ----------
\newcommand{\Lmu}{\mathcal{L}_{\mathrm{BA},\mu}}
\newcommand{\Tfa}{T_{\mathrm{fa}}}
\newcommand{\Q}{\mathbb{Q}}
\newcommand{\N}{\mathbb{N}}

% ---------- Title ----------
\title{Countable Additivity is Not First-Order Axiomatizable}
\author{Patrick S.\ Mahon}
\date{}

\begin{document}
\maketitle

\begin{abstract}
In the first-order language of Boolean algebras equipped with a normalized finitely
additive charge, no first-order theory characterizes those models whose charge extends
to a $\sigma$-additive measure on the generated $\sigma$-algebra.  Equivalently,
countable additivity is not first-order axiomatizable among finitely additive probability
Boolean algebras.  The proof uses a single ultraproduct construction: Dirac charges on
the finite-cofinite algebra are individually $\sigma$-additive, but their ultraproduct
is the finite-cofinite charge, which fails $\sigma$-additivity --- contradicting
\L{}o\'{s}'s theorem if any first-order axiomatization existed.  As an application, the
same obstruction shows that collective exhaustion --- the condition requiring that
whenever a decreasing sequence of observable events has empty intersection, the assigned
charges converge to zero --- cannot be derived from any first-order structural condition,
even though it is the necessary and sufficient condition for $\sigma$-additive extension
in a directed system of Boolean algebras.
\end{abstract}

% ------------------------------------------------------------------
\section{Introduction}
\label{sec:intro}
% ------------------------------------------------------------------

A classical question in model theory is whether a mathematically natural property of
structures can be captured by a first-order theory in an appropriate language.  When
the property is preserved under ultraproducts, \L{}o\'{s}'s theorem ensures it is at
least consistent with first-order axiomatization; when it fails in an ultraproduct of
structures that individually satisfy it, first-order axiomatization is ruled out.

Countable additivity is a fundamental property of probability measures, distinguishing
$\sigma$-additive measures from merely finitely additive charges.  It is well known
that countable additivity cannot be expressed by a single first-order sentence, since
it involves quantification over countable sequences.  The stronger question --- whether
countable additivity is \emph{axiomatizable} by any first-order theory, i.e., whether
there exists a set of first-order sentences whose models are exactly the $\sigma$-additive
ones among finitely additive probability Boolean algebras --- appears not to have been
explicitly settled in the literature in the natural language we consider here.

We settle this question negatively.  In the first-order language $\Lmu$ of Boolean
algebras with a normalized finitely additive charge (Section~\ref{sec:language}), no
first-order theory characterizes those models whose charge extends to a $\sigma$-additive
measure.  The proof is a direct ultraproduct argument using the finite-cofinite algebra
on $\Q$ as the canonical witness (Section~\ref{sec:theorem}).

Section~\ref{sec:ce} draws out an application to collective exhaustion (CE), the
necessary and sufficient admissibility condition for $\sigma$-additive extension in
directed systems of Boolean algebras.  Section~\ref{sec:rep} raises the representation
question of which purely finitely additive charges arise as ultraproducts of
$\sigma$-additive ones, which is strictly harder than non-axiomatizability and is left
open.

% ------------------------------------------------------------------
\section{The language}
\label{sec:language}
% ------------------------------------------------------------------

Let $\Lmu$ be the one-sorted first-order language with:
\begin{itemize}
  \item Boolean operations $\wedge, \vee, \neg, 0, 1$;
  \item for each rational $q \in [0,1] \cap \Q$, unary relation symbols
        $M_{\leq q}(x)$ and $M_{\geq q}(x)$, intended to mean
        $\mu(x) \leq q$ and $\mu(x) \geq q$ respectively.
\end{itemize}

A structure for $\Lmu$ is a Boolean algebra $B$ together with a
$[0,1]$-valued finitely additive probability $\mu$, encoded through the
truth of the predicates $M_{\leq q}$ and $M_{\geq q}$.  Using rational
comparison predicates rather than a function symbol into $[0,1]$ avoids
introducing a second numeric sort.

\begin{definition}[First-order theory of finitely additive probabilities]
Let $\Tfa$ be the $\Lmu$-theory axiomatizing:
\begin{enumerate}[label=(\roman*)]
  \item $B$ is a Boolean algebra;
  \item normalization: $\mu(0) = 0$ and $\mu(1) = 1$;
  \item monotonicity: $x \leq y \Rightarrow \mu(x) \leq \mu(y)$;
  \item finite additivity: for all rationals $r, s$ with $r + s \leq 1$,
        \[
          x \wedge y = 0 \wedge M_{\leq r}(x) \wedge M_{\leq s}(y)
          \;\Rightarrow\; M_{\leq r+s}(x \vee y),
        \]
        and similarly for lower bounds.
\end{enumerate}
\end{definition}

\noindent
The class of Boolean algebras with normalized finitely additive probability is
first-order axiomatizable in $\Lmu$.  Countable additivity, by contrast, requires
quantification over an arbitrary countable sequence $(a_n)_{n \in \N}$ --- not
available in first-order logic --- and so is not a sentence of $\Lmu$.

% ------------------------------------------------------------------
\section{The main theorem}
\label{sec:theorem}
% ------------------------------------------------------------------

\begin{proposition}[$\sigma$-additive extension is not first-order axiomatizable]
\label{prop:not-fo}
There is no first-order $\Lmu$-theory $T \supseteq \Tfa$ such that, for every
$(B, \mu) \models \Tfa$,
\[
  (B, \mu) \models T \quad\Longleftrightarrow\quad
  \mu \text{ extends to a $\sigma$-additive measure on } \sigma(B).
\]
Equivalently, countable additivity is not first-order axiomatizable among
finitely additive probability Boolean algebras.
\end{proposition}

\begin{proof}
Let $E$ be the Boolean algebra of finite-cofinite subsets of $\Q$, fix an
enumeration $q : \N \to \Q$, and let $\delta_n$ be the Dirac charge at
$q(n)$.  Each $(E, \delta_n)$ is a model of $\Tfa$ in which $\delta_n$ is
$\sigma$-additive.  Suppose for contradiction that a theory $T$ as described
exists; then each $(E, \delta_n) \models T$.

Let $\mathcal{U}$ be a nonprincipal ultrafilter on $\N$ extending the cofinite
filter.  By \L{}o\'{s}'s theorem \cite{Los1955}, the ultraproduct $\prod_\mathcal{U}(E,\delta_n)$
also models $T$.  The charge induced on the diagonal copy of $E$ is
\[
  \ell(A) = \lim_\mathcal{U} \delta_n(A) = \lim_\mathcal{U} \mathbf{1}_{q(n) \in A}.
\]
If $A$ is finite, $q(n) \in A$ for only finitely many $n$, so $\ell(A) = 0$.
If $A$ is cofinite, $q(n) \in A$ for cofinitely many $n$, so $\ell(A) = 1$
since $\mathcal{U}$ extends the cofinite filter.  Thus $\ell$ is the
finite-cofinite charge.

The sequence $A_k = \Q \setminus \{q(0), \ldots, q(k)\}$ satisfies
$A_0 \supseteq A_1 \supseteq \cdots$ and $\bigcap_k A_k = \varnothing$, yet
$\ell(A_k) = 1$ for all $k$.  Hence the charge induced on the diagonal copy
of $E$ fails $\sigma$-additivity and does not extend to a $\sigma$-additive
measure on $\sigma(E)$.  Consequently the ultraproduct cannot satisfy any
theory $T$ characterizing $\sigma$-additive extension, contradicting
\L{}o\'{s}'s theorem applied to the assumption that each $(E, \delta_n) \models T$.
\end{proof}

\begin{remark}
For precision: the ultraproduct $\prod_\mathcal{U}(E,\delta_n)$
is an $\Lmu$-structure in its own right; $\ell$ is the charge induced on the
diagonal copy of $E$ inside it.  Identifying the ultraproduct with $(E,\ell)$
is harmless for the contradiction, since the relevant failure of $\sigma$-additivity
is witnessed on the diagonal copy.
\end{remark}

% ------------------------------------------------------------------
\section{Application to collective exhaustion}
\label{sec:ce}
% ------------------------------------------------------------------

Let $(\mathcal{E}_i, \ell_i)_{i \in \iota}$ be a directed system of Boolean charge
spaces: each $\mathcal{E}_i$ is a Boolean algebra and $\ell_i : \mathcal{E}_i \to [0,1]$
is a normalized finitely additive charge, with compatible connecting maps.
\emph{Collective exhaustion} (CE) is the following condition on the induced charge
$\ell$ on the direct limit algebra $\mathcal{E} = \varinjlim \mathcal{E}_i$: whenever
$E_1 \supseteq E_2 \supseteq \cdots$ in $\mathcal{E}$ with $\bigcap_n E_n = \varnothing$,
we have $\ell(E_n) \to 0$.

CE is a valuation-layer condition: it constrains how charges behave across the directed
system, not the Boolean-algebraic structure of the index set.  It is the necessary and
sufficient condition under which a compatible finitely additive family extends to a
$\sigma$-additive measure on the observable $\sigma$-algebra generated by the system.

Proposition~\ref{prop:not-fo} implies immediately:

\begin{corollary}[CE is not first-order]
\label{cor:ce-not-fo}
In the first-order language $\Lmu$ of Boolean algebras with normalized finitely
additive charge, no first-order $\Lmu$-theory characterizes those models satisfying CE.
Equivalently, no accumulation of first-order conditions on the Boolean-algebraic
structure can imply collective exhaustion.
\end{corollary}

\noindent
The finite-cofinite charge of Proposition~\ref{prop:not-fo} is the canonical witness:
it satisfies every first-order structural condition yet assigns persistent unit mass to
the decreasing sequence $A_k = \Q \setminus \{q(0), \ldots, q(k)\}$, whose intersection
is empty.  CE fails precisely because the mass is not exhausted.

Foundational disputes about probability are not terminological.  They concern which
extra structure licenses the passage from finite coherence to probability, and
Corollary~\ref{cor:ce-not-fo} shows that some such structure is logically unavoidable:
no first-order condition can close the gap.  For a full development of CE as the
admissibility condition for $\sigma$-additive extension in a directed observational
system, see \cite{mahon2025}.

% ------------------------------------------------------------------
\section{The representation question}
\label{sec:rep}
% ------------------------------------------------------------------

Proposition~\ref{prop:not-fo} shows that \emph{one} purely finitely additive charge
arises as an ultraproduct of $\sigma$-additive ones --- that single witness suffices
for the non-axiomatizability conclusion.  A stronger question is:

\begin{quote}
\emph{Which purely finitely additive charges arise as ultraproducts or ultralimits
of $\sigma$-additive probabilities?}
\end{quote}

This is a representation question, and it is strictly harder than non-axiomatizability.
It may hold only after enlarging the notion of approximation, for example by
allowing nets or ultralimits more general than literal ultraproducts.  The
Yosida--Hewitt decomposition \cite{YosidaHewitt1952} identifies the purely
finitely additive part $\ell_p$ of any finitely additive charge; whether $\ell_p$
can always be represented as an ultralimit is not settled here and is left as an
open question.

\bibliographystyle{plain}
\bibliography{../references}

\end{document}
